\section{Introduction}

Autonomous underwater vehicles (AUVs) play an increasingly vital role in ocean exploration, subsea inspection, and environmental monitoring \cite{fossen2011handbook}. During trajectory tracking missions, ocean currents represent the dominant disturbance source affecting control accuracy. Accurate estimation of current velocity for feedforward compensation is therefore critical to achieving robust AUV control performance \cite{borhaug2007model}.

Existing AUV current estimation methods can be broadly categorized into two paradigms. The first comprises kinematic current observers \cite{liang2018three,he2024three}, which estimate inertial-frame current components through integration of position or velocity errors. These methods do not rely on hydrodynamic model accuracy and are thus inherently robust. However, their convergence is limited by the accumulation rate of position errors, typically requiring tens of seconds, making them slow to respond to abrupt current changes. The second paradigm encompasses extended state observer (ESO) based lumped disturbance estimation \cite{xie2020extended}, which treats ocean currents, model uncertainty, and unmodeled dynamics as a single combined disturbance signal. While responsive, ESO-based methods produce a compensation signal rather than a physically interpretable current estimate, as they cannot separate physical current effects from model errors.

A practical challenge that cuts across both paradigms is \emph{hydrodynamic model mismatch}. Kinematic observers are naturally immune, as they do not use a dynamics model. ESO-based methods absorb model error into the lumped disturbance indiscriminately. However, when researchers attempt to leverage dynamics models for direct physical current estimation---as in the high-gain observer with online recursive least squares identification by \cite{kim2018estimating}, or the nonlinear Luenberger observer by \cite{borhaug2007model} which assumes a perfectly known dynamic model---model mismatch can directly contaminate the current estimate. This leads to the central question of this paper: \textbf{Can we design a current observer that exploits the information provided by a dynamics model while remaining robust to model mismatch?}

We propose an affirmative answer. Our approach is motivated by an engineering setting in which the XHY AUV possesses a CFD (RANS $k$-$\omega$ SST) pre-calibrated hydrodynamic model with per-degree-of-freedom drag coefficient goodness-of-fit $R^2>0.99$, independently validated through still-water towing tests and free-running trials. This provides not merely a nominal model, but a \emph{deterministic prior} with a quantifiable accuracy bound.

Building on this, we present the CFD-Prior Current Observer (PC-RCO), a prediction-correction architecture with two key innovations:

\begin{enumerate}[itemsep=0pt,topsep=4pt]
    \item \textbf{Continuous Gauss-Markov (GM) regularization}: Rather than relying on a binary excitation gate that may fail to close under feedback control, PC-RCO applies a weak GM decay toward the prior mean at \emph{every} time step. During high-excitation phases, the gradient-driven update dominates; during low-excitation or high-noise phases, the GM regularization gradually pulls the estimate toward the prior, preventing drift while maintaining bounded updates. This continuous blending eliminates the need for threshold-based gating, which we show cannot reliably distinguish excitation levels in feedback-controlled AUVs.
    \item \textbf{Adaptive update-rate clamping with online noise estimation}: PC-RCO maintains an exponential moving average (EMA) estimate of the instantaneous DVL noise level from the velocity prediction residual. The per-step update bound ($\Delta c_{\max}$) is dynamically scaled based on this noise estimate---tight during low-noise conditions for smooth convergence, relaxed during high-noise conditions to maintain convergence through noisy measurements. An absolute upper bound ($0.20$ m/s) prevents unbounded updates under any conditions.
\end{enumerate}

We conduct systematic numerical validation of PC-RCO on the XHY AUV simulation platform, covering 6 CFD drag model mismatch levels (0--50\%), 4 baseline methods spanning 3 methodological categories, and controlled DVL noise conditions. Results demonstrate: (1) PC-RCO degrades only 42\% under 30\% CFD drag perturbation versus 77\% for the EKF \cite{long2021trajectory}, while achieving 3$\times$ better accuracy than the kinematic observer at zero mismatch (RMSE$_{V_c}=0.064$ vs.\ $0.192$); (2) continuous GM regularization alone contributes 57--218\% improvement in estimation accuracy under model mismatch and high-noise conditions, as confirmed by ablation; (3) adaptive clamping reduces the maximum per-step estimation jump by 34$\times$ under high DVL noise ($\sigma=0.10$ m/s) compared to an unclamped variant, preventing physically impossible transient estimates (3 m/s single-step jumps) without sacrificing steady-state accuracy.

The remainder of this paper is organized as follows. Section~2 formulates the AUV dynamics model, current parameterization, and CFD uncertainty characterization. Section~3 details the PC-RCO observer design. Section~4 presents systematic numerical validation. Section~5 discusses applicability, limitations, and deployment considerations. Section~6 concludes the paper.
